% !TEX program = xelatex
% The research corpus (ci/research.json): polyglossia with Arabic,
% Hebrew (right to left, bidi), Greek and Russian, on TeX Live's fonts.
\documentclass[11pt,a4paper]{article}
\usepackage{fontspec}
\usepackage{polyglossia}
\setdefaultlanguage{english}
\setotherlanguages{arabic,hebrew,greek,russian}
\setmainfont{Libertinus Serif}
\setsansfont{Libertinus Sans}
\newfontfamily\arabicfont[Script=Arabic]{Amiri}
\newfontfamily\hebrewfont[Script=Hebrew]{David CLM}
\newfontfamily\hebrewfontsf[Script=Hebrew]{Frank Ruehl CLM}
\title{Four Scripts in One Paper}
\author{K. Linguist}
\date{2026}
\begin{document}
\maketitle
\section{Greek}
\begin{greek}
Ἐν ἀρχῇ ἦν ὁ λόγος, καὶ ὁ λόγος ἦν πρὸς τὸν θεόν. Η γλώσσα είναι
ζωντανή, και οι λέξεις αλλάζουν με τον χρόνο· η συλλαβοποίηση ακολουθεί
τους κανόνες της ελληνικής γλώσσας.
\end{greek}
\section{Russian}
\begin{russian}
Все счастливые семьи похожи друг на друга, каждая несчастливая семья
несчастлива по-своему. Перенос слов следует правилам русского языка.
\end{russian}
\section{Arabic}
\begin{arabic}
بسم الله الرحمن الرحيم. اللغة العربية من أقدم اللغات السامية، وتكتب من
اليمين إلى اليسار، والأرقام ١٢٣٤٥٦٧٨٩٠.
\end{arabic}
\section{Hebrew}
\begin{hebrew}
בראשית ברא אלהים את השמים ואת הארץ. העברית נכתבת מימין לשמאל.
\end{hebrew}
\section{Mixed}
English with an inline \textgreek{λόγος}, \textrussian{слово},
\textarabic{كلمة} and \texthebrew{מילה}, then back to English.
\end{document}
